% Provenance: copied from the design workspace CambridgePQ/recursion.fast, `report/sections/03-background.tex`, on 2026-10-19. Read-only copy; the source is authoritative.
% 03-background -- written by W1 (background for cryptographers and formal-methods readers)
% Unicode in listings: listings reorders non-ASCII glyphs under LuaLaTeX unless they are
% declared as literates (they would otherwise be emitted out of order, e.g. after a parenthesis).
\section{Background}\label{sec:background}

This section fixes notation and records the state of the art on which the competition design
rests. It has two audiences. Cryptographers should find
\cref{subsec:bg-lean} (what a Lean proof does and does not certify) self-contained. Formal-methods
readers should find the same for \cref{subsec:bg-r1cs,subsec:bg-iop,subsec:bg-bcs,subsec:bg-proximity,subsec:bg-binary,subsec:bg-costs,subsec:bg-recursion}.
Claims marked \emph{(unverified)} could not be confirmed from a primary source; repository facts
are cited as \repo{repo@commit} with a path, from the research notes behind this report.

Throughout, $\F$ is a finite field, $\lambda$ is a security parameter in bits, and
$[n]=\{1,\dots,n\}$. A probability bound is \emph{concrete} if it is stated as a function of the
parameters, without asymptotic notation. All adversaries in information-theoretic statements are
computationally unbounded.

%-----------------------------------------------------------------------------------------------
\subsection{R1CS, indexed relations and holography}\label{subsec:bg-r1cs}

\begin{definition}[R1CS]\label{def:bg-r1cs}
Fix sizes $m$ (constraints), $n_x$ (public inputs) and $n_w$ (witness variables), and let
$N = 1+n_x+n_w$. An \emph{R1CS index} is a triple $i=(A,B,C)$ of matrices in $\F^{m\times N}$.
An \emph{instance} is $x\in\F^{n_x}$ and a \emph{witness} is $w\in\F^{n_w}$. With
$z=(1,x,w)\in\F^N$,
\[
  (i,x,w)\in\Rrcs \iff (A z)\circ(B z) = C z ,
\]
where $\circ$ is the entrywise (Hadamard) product. The \emph{sparsity} is
$\mathsf{nnz}(i) = |\mathrm{supp}\,A|+|\mathrm{supp}\,B|+|\mathrm{supp}\,C|$, the number of nonzero
entries (some papers use the maximum of the three counts instead; the two differ by a factor of at
most $3$).
\end{definition}

The leading constant $1$ in $z$ lets affine constraints be written; ArkLib's
\lean{R1CS.relation} uses $z=(x,w)$ without it, so the competition restates the relation
locally.\footnote{\repo{ArkLib@db7005f}, \path{ArkLib/ProofSystem/ConstraintSystem/R1CS.lean:64}
(\lean{relation}); its field comments for $m$ and $n$ are also swapped (lines 32--33).}
recursion.fast instantiates $\F$ with the binary field $K=\mathrm{GF}(2^{64})$
(\cref{subsec:bg-binary}). R1CS is NP-complete, and a RISC-V execution trace, a hash chain or a SNARK verifier can all be
arithmetised into it; the matrices are then very sparse ($\mathsf{nnz}=O(m+N)$).

\paragraph{Indexed relations.} An \emph{indexed relation} $\R$ is a set of triples $(i,x,w)$:
the index $i$ is large and fixed in advance (here the matrices), the instance $x$ is small and
varies from proof to proof. Its language is $L(\R)=\{(i,x) : \exists w,\ (i,x,w)\in\R\}$
\cite{CHMMVW20-Marlin,COS20}, \cite[ch.~32]{CY24-Book}. Soundness for an indexed relation
quantifies over \emph{adversarially chosen} indices: the cheating prover outputs $(i,x)$ (and
auxiliary data) itself \cite[Def.~32.4.2]{CY24-Book}.

\paragraph{Why succinct verification of arbitrary R1CS needs preprocessing.} An arbitrary
(non-uniform) index has description length $\Theta(\mathsf{nnz}\cdot\log(mN|\F|))$ bits. A
verifier that receives $i$ as input therefore runs in time $\Omega(\mathsf{nnz})$ just to read it,
and checking $(Az)\circ(Bz)=Cz$ directly costs $\Theta(\mathsf{nnz})$ field operations. Succinct
verification is possible only if the cost of reading $i$ is paid once, by an \emph{indexer}.

\begin{definition}[Holographic proof system]\label{def:bg-holographic}
A holographic (preprocessing) proof system for $\R$ is a triple $(\Ind,\Prv,\Vrf)$. The
deterministic indexer computes, once per index, an encoding $\Ind(i)$ (the \emph{index oracle})
and a short verification key $\vk_i$. The prover receives $(i,x,w)$. The verifier receives
$(\vk_i,x)$ and \emph{oracle access} to $\Ind(i)$, and runs in time polylogarithmic in
$|i|$. In the compiled argument, $\vk_i$ contains a Merkle root of $\Ind(i)$, and the verifier's
queries to $\Ind(i)$ become authenticated openings \cite[ch.~32]{CY24-Book}.
\end{definition}

\paragraph{Spartan, Marlin, Fractal and COS.} \emph{Spartan} \cite{Set20-Spartan} reduces R1CS
satisfiability to a sumcheck over the multilinear extensions $\tilde A,\tilde B,\tilde C$ and
$\tilde z$. At the end the verifier needs $\tilde A(r_x,r_y)$, $\tilde B(r_x,r_y)$,
$\tilde C(r_x,r_y)$ at a random point. SPARK, Spartan's sparse-polynomial compiler, lets the
prover prove these evaluations against a commitment to the matrices computed in preprocessing
(Spartan's ``computation commitment''), so the verifier costs $O(\log^2 n)$ field operations plus
polynomial-commitment checks. \emph{Marlin} \cite{CHMMVW20-Marlin} introduced holographic
polynomial IOPs for R1CS, in which the indexer encodes each matrix by low-degree polynomials
(row, column, value) that the verifier queries. \emph{Fractal} \cite{COS20} gave holographic IOPs
with Reed--Solomon encodings and the COS transformation, the holographic analogue of BCS
(\cref{subsec:bg-bcs}). The result is transparent, plausibly post-quantum and recursive, and
\cite[ch.~32]{CY24-Book} formalises it. Pairing Spartan's information-theoretic core with a
hash-based polynomial commitment (WHIR or BaseFold, \cref{subsec:bg-costs}) yields a
hash-based R1CS verifier; STIR estimates about $220$\,KiB for an $n=2^{24}$ R1CS at $128$ bits and
$\rho=1/2$ \cite{ACFY24-STIR}.

%-----------------------------------------------------------------------------------------------
\subsection{Interactive oracle proofs and round-by-round knowledge soundness}\label{subsec:bg-iop}

\paragraph{IOPs.} An \emph{interactive oracle proof} (IOP)
\cite{BCS16,RRR16} with $k$ rounds runs as follows. In round $j\in[k]$ the prover sends an
\emph{oracle} $\Pi_j\in\Sigma_j^{\ell_j}$, a string the verifier may later read at chosen
positions. The verifier then replies with a challenge $\rho_j$. After the last round the
verifier queries the oracles (and $\Ind(i)$, if holographic) and accepts or rejects. The IOP is
\emph{public-coin} if every $\rho_j$ is uniform on a finite set $C_j$ and independent of
everything else, so that the verifier's randomness is its messages; only public-coin IOPs can be
compiled by Fiat--Shamir. We write $\mathrm{tr}_{<j}$ for the partial transcript before round $j$,
$\ell=\sum_j\ell_j$ for the total oracle length, and $q$ for the number of queries. An \emph{IOP of
proximity} (IOPP) instead gives the verifier oracle access to its input $f$ and asks it to reject
when $f$ is $\delta$-far from a code (\cref{subsec:bg-proximity}). FRI, STIR and WHIR are IOPPs.

Fix an indexed relation $\R$ and a public-coin holographic IOP $(\Ind,\Prv,\Vrf)$. Write
$\langle \tilde\Prv,\Vrf\rangle(i,x)\in\{0,1\}$ for the verifier's decision when it interacts
with a (possibly malicious) prover $\tilde\Prv$, the probability being over the challenges.

\begin{definition}[Completeness, soundness, knowledge soundness]\label{def:bg-ks}
The IOP has
\begin{itemize}[nosep]
  \item \emph{perfect completeness} if $\Pr[\langle\Prv(i,x,w),\Vrf\rangle(i,x)=1]=1$ for every
    $(i,x,w)\in\R$;
  \item \emph{soundness error} $\eps$ if $\Pr[\langle\tilde\Prv,\Vrf\rangle(i,x)=1]\le\eps$ for
    every $(i,x)\notin L(\R)$ and every $\tilde\Prv$;
  \item \emph{(straight-line) knowledge error} $\kappa$ if there is an extractor $\Ext$ such that for
    every $(i,x)$ and every $\tilde\Prv$,
    \[ \Pr\bigl[\langle\tilde\Prv,\Vrf\rangle(i,x)=1 \ \wedge\ (i,x,\Ext(i,x,\Pi_1,\dots,\Pi_k))\notin\R\bigr]\le\kappa, \]
    where $\Ext$ sees the prover's oracles but not its code (it is \emph{straight-line}, i.e.\ it
    does not rewind). Cryptographic definitions require $\Ext$ to run in polynomial time.
\end{itemize}
\end{definition}

Completeness is not a formality in an adversarial competition: a verifier that rejects
everything is trivially sound, and only completeness excludes it.

\paragraph{Round-by-round soundness.} Soundness bounds the total probability of acceptance;
Fiat--Shamir needs a stronger, per-round property. Round-by-round (RBR) soundness, introduced by
Canetti et al.\ \cite{CCHLRRW19}, requires that a false statement stay ``doomed'' after each
round except with small probability over that round's challenge.

\begin{definition}[State function; RBR soundness {\cite{CCHLRRW19}, \cite[Defs.~31.1.1--31.1.2]{CY24-Book}}]\label{def:bg-rbr}
A \emph{state function} is a deterministic (possibly inefficient) predicate
$\mathsf{St}((i,x),\mathrm{tr})\in\{0,1\}$ on partial transcripts such that, for $(i,x)\notin L(\R)$:
(i) $\mathsf{St}((i,x),\emptyset)=0$; (ii) prover moves cannot leave the doomed state:
$\mathsf{St}(\mathrm{tr})=0\Rightarrow\mathsf{St}(\mathrm{tr}\|\Pi)=0$ for every $\Pi$; and
(iii) if $\mathsf{St}(\mathrm{tr})=0$ on a full transcript then $\Vrf$ rejects. The IOP has
\emph{RBR soundness error} $(\eps_1,\dots,\eps_k)$ if some state function satisfies, for every
$(i,x)\notin L(\R)$, every $j$ and every prefix $\mathrm{tr}\|\Pi_j$ with
$\mathsf{St}(\mathrm{tr}\|\Pi_j)=0$,
\[ \Pr_{\rho_j\leftarrow C_j}\bigl[\mathsf{St}(\mathrm{tr}\|\Pi_j\|\rho_j)=1\bigr]\le\eps_j . \]
\end{definition}

Only verifier randomness can flip a state from $0$ to $1$, so a union bound gives soundness error
at most $\sum_j\eps_j\le k\cdot\max_j\eps_j$ \cite[Claim~31.1.3]{CY24-Book}.

\paragraph{RBR knowledge soundness: three variants.} The knowledge analogue
\cite{CMS19,COS20} comes in (at least) three variants that differ in what the extractor sees:
CMS19's, whose extractor sees the transcript prefix \cite[Def.~8.5]{CMS19}; the book's, whose
single extractor sees all prover strings of the continued execution \cite[Def.~31.1.6]{CY24-Book}
(\cref{def:bg-rbrks-cy}); and WARP's per-round backward extractor \cite[Defs.~4.1--4.2]{BCFW25-WARP}
(\cref{def:bg-rbrks-warp}). The competition must pin exactly one of them (\cref{sec:security}).

\begin{definition}[RBR-KS, global extractor; paraphrasing {\cite[Def.~31.1.6]{CY24-Book}}]\label{def:bg-rbrks-cy}
There are a state function $\mathsf{St}$ and a polynomial-time extractor $\Ext$ that sees only
prover strings (hence is straight-line) such that for every $(i,x)$, every malicious prover
$\tilde\Prv$, every $j$ and every prefix with $\mathsf{St}(\mathrm{tr}\|\Pi_j)=0$,
\[ \Pr_{\rho_j,\ \tilde\Prv}\bigl[\mathsf{St}(\mathrm{tr}\|\Pi_j\|\rho_j)=1 \ \wedge\
   (i,x,\Ext(i,x,\vec\Pi))\notin\R\bigr]\le\kappa_j , \]
where the probability is over $\rho_j$ and the continuation $(\Pi_{j+1},\dots,\Pi_k)$ produced by
$\tilde\Prv$, and $\vec\Pi$ is the list of \emph{all} prover oracles of that execution, including
those sent after round $j$.
\end{definition}

The second variant replaces the single extractor by a chain of per-round maps that pull a witness
backwards through the transcript. It is due to WARP \cite[Defs.~4.1--4.2, informally
2.6--2.7]{BCFW25-WARP} and is the
form ArkLib implements (\lean{KnowledgeStateFunction} with \lean{Extractor.RoundByRound}).
Position the $n=2k$ messages in order (prover message, then challenge) and index stages by
$m\in\{0,\dots,n\}$.

\begin{definition}[RBR-KS, per-round backward extractor {\cite[Def.~4.2, informally Def.~2.7]{BCFW25-WARP}}]\label{def:bg-rbrks-warp}
A \emph{knowledge state function} is a family of predicates $S_m(s,\mathrm{tr},w)$ on statements
$s=(i,x)$, stage-$m$ transcripts and intermediate witnesses $w\in W_m$ (with $W_0$ the witness
type of $\R$). It comes with extractor maps $E_m:\ (s,\mathrm{tr}\|a,w)\mapsto w'$ from $W_{m+1}$
to $W_m$ and $E_{\mathrm{out}}$ into $W_n$, satisfying three laws:
\begin{itemize}[nosep]
  \item (empty) $S_0(s,\emptyset,w)\iff (s,w)\in\R$;
  \item (prover moves) if message $m$ is the prover's, then
    $S_{m+1}(s,\mathrm{tr}\|\Pi,w)\Rightarrow S_m(s,\mathrm{tr},E_m(s,\mathrm{tr}\|\Pi,w))$;
  \item (full) if $\Vrf(s,\mathrm{tr})=\mathsf{some}\ s'$ and $(s',w')\in\R_{\mathrm{out}}$, then
    $S_n(s,\mathrm{tr},E_{\mathrm{out}}(s,\mathrm{tr},w'))$; for a proof system,
    $\R_{\mathrm{out}}$ is trivial (\cref{def:bg-rok}).
\end{itemize}
For a challenge position $m$ and prefix $\mathrm{tr}$, a challenge $c\in C_m$ is \emph{bad} if
$\exists w:\ \neg S_m(s,\mathrm{tr},E_m(s,\mathrm{tr}\|c,w))\wedge S_{m+1}(s,\mathrm{tr}\|c,w)$.
The verifier is \emph{(worst-case) RBR knowledge sound} with errors $(\eps_m)$ if, for every $s$,
every challenge position $m$ and every prefix $\mathrm{tr}$,
\[ \#\{c\in C_m : c \text{ bad}\}\ \le\ \eps_m\cdot|C_m| . \]
\end{definition}

\begin{remark}[How the variants differ]\label[remark]{rem:bg-variants}
(1) \emph{Extractor shape.} \Cref{def:bg-rbrks-cy} has one extractor that sees every prover
string. \Cref{def:bg-rbrks-warp} has one local map per round, and extraction only needs to
\emph{fail} at a challenge. WARP shows that its notion is implied by the CMS19-style notion
(adapted to reductions) \cite[App.~C]{BCFW25-WARP} and states that the relaxation is strict, which
is what admits codes not known to have an efficient \emph{error-tolerant} decoder; it does not
compare its notion with the book's. ArkLib further
generalises the definition in two documented ways: the maps may transform the witness across
prover moves, and the full-transcript law is one-directional.\footnote{\repo{ArkLib@db7005f},
\path{ArkLib/OracleReduction/Security/RoundByRound.lean:155-164} (docstring: deviations from
ABF26 Def.~A.5, in ArkLib's numbering).} (2) \emph{Worst case versus average.} \Cref{def:bg-rbrks-warp} quantifies over
every prefix, which matches the literature. ArkLib's default \lean{rbrKnowledgeSoundness} instead
lets the prover sample the prefix and bounds the mixture, which is weaker; the worst-case form
\lean{rbrKnowledgeSoundnessWorstCase} exists, and the implication worst-case $\Rightarrow$
averaged is proved.\footnote{\repo{ArkLib@db7005f},
\path{ArkLib/OracleReduction/Security/RoundByRound.lean:539} (definition) and \texttt{:609}
(implication).} (3) \emph{Efficiency.} Cryptographic definitions require a polynomial-time
extractor; ArkLib's extractor maps are arbitrary Lean functions (\cref{subsec:bg-lean}).
\end{remark}

\paragraph{State-restoration soundness.} The property that Fiat--Shamir and BCS actually consume
is \emph{state-restoration} (SR) soundness \cite{BCS16}, \cite[Defs.~23.2.2--23.2.3]{CY24-Book}.
In the SR game the malicious prover talks to a verifier that it may \emph{rewind}: at each of at
most $t$ moves it picks any previously reached verifier state (a partial transcript), appends a
prover message, and receives a fresh challenge. It wins if some complete transcript in the
resulting tree is accepting while $(i,x)\notin L(\R)$ (soundness, error $\eps^{\mathrm{sr}}(t)$),
or is accepting while the extractor, given the tree, fails to output a witness (knowledge, error
$\kappa^{\mathrm{sr}}(t)$). This models a Fiat--Shamir prover, which can recompute a hash on any
modified prefix.

\begin{theorem}[Relations between the notions]\label[theorem]{thm:bg-rbr-sr}
Let $\eps_{\mathrm{rbr}}=\max_j\eps_j$ and $\kappa_{\mathrm{rbr}}=\max_j\kappa_j$.
\begin{enumerate}[nosep]
  \item RBR-KS with errors $(\kappa_j)$ implies RBR soundness with the same errors.
  \item RBR soundness implies soundness error $\sum_j\eps_j$, and RBR-KS implies straight-line
    knowledge error $\sum_j\kappa_j$ \cite[Claim~31.1.3]{CY24-Book}.
  \item $\eps^{\mathrm{sr}}(t)\le(t+k)\cdot\eps_{\mathrm{rbr}}$ \cite[Thm.~31.2.1]{CY24-Book}, and
    $\kappa^{\mathrm{sr}}(t)\le(t+k)\cdot\kappa_{\mathrm{rbr}}$ with a straight-line SR extractor
    for the book's notion \cite[Thm.~31.3.1]{CY24-Book}. For the backward-extractor variant, the
    same $(t+k)\cdot\kappa_{\mathrm{rbr}}$ bound is \cite[Thm.~2.8, formally Thm.~B.4 with
    Constr.~B.5]{BCFW25-WARP}, with extraction time $\sum_i\mathit{et}_i$.
  \item For interactive proofs, RBR soundness and SR soundness are qualitatively (asymptotically) equivalent, and
    neither follows from Fiat--Shamir soundness in the random-oracle model \cite{Hol19}.
\end{enumerate}
\end{theorem}

Item~1 is proved in ArkLib.\footnote{\repo{ArkLib@db7005f},
\path{ArkLib/OracleReduction/Security/Implications.lean:87}.} The implications to plain
soundness and knowledge soundness in item~2, and all SR implications in ArkLib's legacy layer, are
currently \lean{sorry}. ArkLib's newer native \lean{Interaction} layer proves an SR knowledge bound
$(Q+\#\text{rounds})\cdot\eps$, sorry-free, for a fixed-round presentation with finite message and
challenge types. It is not yet connected to the legacy definitions
(\cref{subsec:bg-lean}).\footnote{\repo{ArkLib@db7005f},
\path{ArkLib/Interaction/Oracle/Security/StateRestorationGame.lean:102}
(\lean{stateRestoration\_knowledge\_soundness}).}

The alternative route to Fiat--Shamir security is \emph{special soundness}. For multi-round
special-sound protocols, Attema, Fehr and Klo\ss{} show that the Fiat--Shamir knowledge error
grows only linearly in the number of random-oracle queries \cite{AFK22}. Block et al.\ relate the
generalised notions for IOPs: generalised special soundness implies generalised RBR soundness,
and generalised RBR-KS implies generalised special soundness, but not conversely
\cite{BGTZ23}. For hash-based IOPs, RBR (knowledge) soundness is the notion that has actually been
proven: for FRI's Fiat--Shamir security \cite{BGKTTZ23-FSFRI}, for STIR (RBR soundness)
\cite{ACFY24-STIR} and for WHIR (RBR-KS) \cite{ACFY24-WHIR}.

%-----------------------------------------------------------------------------------------------
\subsection{From IOPs to SNARKs: BCS and Fiat--Shamir}\label{subsec:bg-bcs}

\paragraph{The BCS transformation.} Ben-Sasson, Chiesa and Spooner \cite{BCS16} compile a
public-coin IOP into a non-interactive argument in the random-oracle model (ROM) as follows. The
prover commits to each oracle $\Pi_j$ with a Merkle tree, whose leaves and nodes are hashed with
the random oracle $H$. It derives $\rho_j$ by hashing the previous state and the new root
$\mathrm{rt}_j$ (with a salt). It derives the query positions from the final state, and it sends
the queried symbols with their authentication paths. The verifier recomputes every challenge and
checks every path. Let $t$ be the adversary's query budget to $H$, $\lambda$ the output length of
$H$ in bits, $k$ the round count and $\ell$ the total oracle length.

\begin{theorem}[BCS, concrete bounds {\cite[Thms.~25.2.1 and 26.1.1 (straight-line clause)]{CY24-Book}}]\label[theorem]{thm:bg-bcs}
For $t\ge 2(\log\ell+1)\cdot\ell$, with $\ell=\sum_j\ell_j$ (the book's condition under which the
Merkle and collision terms combine), the BCS argument has soundness and straight-line knowledge errors
\[
  \eps_{\mathrm{BCS}}(t)\le (t+k)\cdot\eps_{\mathrm{rbr}} + \frac{3.5\,t^2}{2^{\lambda}},\qquad
  \kappa_{\mathrm{BCS}}(t)\le (t+k)\cdot\kappa_{\mathrm{rbr}} + \frac{3.5\,t^2}{2^{\lambda}} .
\]
The extractor runs in time $\mathit{et}^{\mathrm{sr}}+\mathit{et}^{\mathrm{mm}}_{\mathrm{MT}}
+O(r_{\max}\cdot k\cdot t\log t)$: the SR extractor's time, the Merkle extractor's time, and a
term with $r_{\max}=\max_j\log|C_j|$ the largest challenge length. (The book's main knowledge
statements are for rewinding extractors; straight-line extraction is their \enquote{moreover}
clause.) For the holographic COS variant, which commits to $\Ind(i)$ once and places the root in
$\vk_i$, the bound gains a term $4\ell_0^2/2^\lambda$ for an index oracle of length $\ell_0$; with
$\ell:=\ell_0+\sum_j\ell_j$ and $t\ge2\ell_0$ it is $(t+k)\cdot\eps_{\mathrm{rbr}}+4.5\,t^2/2^\lambda$
(soundness \cite[Thm.~32.8.9]{CY24-Book}; knowledge, straight-line clause, \cite[Thm.~32.8.10]{CY24-Book}).
\end{theorem}

\paragraph{Reading the bound.} The factor $(t+k)$ is the price of state restoration. A Lean bound
of $2^{-\lambda_r}$ per round therefore yields only about $2^{-(\lambda_r-\tau)}$ against an
adversary making $t=2^\tau$ hash queries. For example, with
$\eps_{\mathrm{rbr}}\le 2^{-128}$, $\lambda=256$ and $2k\le t\le 2^{124}$ (with $t$ replaced by
$\max(t,t^*)$ in the hash term below the side condition's threshold $t^*$, \cref{subsec:sec-cap}),
we get $\kappa(t)\le 1.5\,t/2^{128}+4.5\,t^2/2^{256}\le t/2^{127}$, which is sig.golf's $Q/2^{127}$
convention (\cref{sec:prior}). With $k\le t$ alone, the first term already equals $t/2^{127}$ and
the bound exceeds it by the (negligible) collision term. Grinding (proof-of-work) adds bits to
individual rounds \cite{ethSTARK21}. Deployed systems implement
Fiat--Shamir with a duplex sponge rather than independent random oracles; Chiesa and Orr\`u
analyse that setting in the ideal-permutation model \cite{CO25-DuplexFS}.

\begin{figure}[t]
\centering
\begin{tikzpicture}[
  box/.style={draw,rounded corners=2pt,align=center,font=\small,minimum height=1.25cm,text width=2.6cm,inner sep=3pt},
  lean/.style={box,fill=green!8,draw=green!45!black},
  cited/.style={box,fill=blue!5,draw=blue!45!black},
  heur/.style={box,fill=orange!8,draw=orange!60!black},
  arr/.style={-{Latex[length=2mm]},thick},
  lab/.style={font=\scriptsize,align=center}]
\node[lean]  (iop) {Holographic IOP\\ completeness +\\ RBR-KS $\eps_{\mathrm{rbr}}$};
\node[cited, right=1.55cm of iop] (sr) {SR knowledge\\ soundness\\ $(t+k)\,\eps_{\mathrm{rbr}}$};
\node[cited, right=1.55cm of sr] (bcs) {BCS/COS argument\\ in the ROM\\ $+\,4.5\,t^2/2^{\lambda}$};
\node[heur, right=1.55cm of bcs] (fs) {Deployed NARG\\ concrete hash\\ (heuristic)};
\draw[arr] (iop) -- node[lab,above]{CY 31.2.1 (sound)}node[lab,below]{WARP B.4 (knowl.)} (sr);
\draw[arr] (sr) -- node[lab,above]{CY 32.8.9 (sound)}node[lab,below]{CY 32.8.10 (knowl.)} (bcs);
\draw[arr] (bcs) -- node[lab,above]{$H\mapsto$ hash}node[lab,below]{binding rule} (fs);
\node[lab,below=0.15cm of iop] {proved in Lean\\ by submitters};
\node[lab,below=0.15cm of sr] {sorry-free only in\\ ArkLib's native layer};
\node[lab,below=0.15cm of bcs] {cited meta-theorem\\ (no Lean BCS)};
\node[lab,below=0.15cm of fs] {KRS25, Fen26:\\ outside any IOP theorem};
\end{tikzpicture}
\caption{The security chain from an IOP to a deployed succinct argument, and which link is
machine-checked. Only the leftmost box is a Lean obligation in recursion.fast; the arrows are
theorems of the cited works; the last arrow is a heuristic that contract rules must protect.}
\label{fig:bg-pipeline}
\end{figure}

\paragraph{Fiat--Shamir pitfalls outside the IOP model.} Three classes of failure sit outside
\cref{thm:bg-bcs} because they concern how the oracle is instantiated or what is hashed.
\begin{enumerate}[nosep,label=(\roman*)]
\item \emph{Weak Fiat--Shamir.} Deriving challenges without hashing the statement lets a prover
  choose the statement after the challenges \cite{BPW12-WeakFS}. In the indexed setting the
  statement is $(i,x)$, so $\vk_i$ \emph{and} $x$ must be absorbed.
\item \emph{Hash inside the relation.} Khovratovich, Rothblum and Soukhanov give, for every
  concrete Fiat--Shamir hash, a circuit for which Fiat--Shamir-compiled GKR is not adaptively
  sound, even under the \emph{strong} transform that absorbs the statement (the circuit
  description). The circuit evaluates the hash on a transcript-predicting input \cite{KRS25}. Recursion
  meets this precondition by construction, because the proved circuit contains the verifier's
  hash.
\item \emph{Generated statements.} Fenzi extends this to R1CS instances \emph{generated by running
  an adversarial program}: when the conversion is expressive enough, the author of the program can
  break Fiat--Shamir soundness of any protocol in which an accepting transcript can be prepared
  before the statement is bound. Variants of Spartan and Aurora \cite{BCRSVW19-Aurora} are
  affected \cite{Fen26-GenR1CS}. The proven mitigation is to derive the first challenge from the
  \emph{generated statement}, meaning the R1CS (or its commitment) and the public input; soundness
  then reduces to that of the underlying protocol for the non-generated relation. (We read this
  paper only at the level of its abstract and the R7 summary.)
\end{enumerate}
In recursion.fast the R1CS is generated from programs, and the next layer's circuit contains the
hash. The design therefore makes the initial absorption a property of the trusted harness, not of
the submission: the transcript starts from
$H(\mathit{domSep}\,\|\,\mathit{contractId}\,\|\,\mathit{shapeHash}\,\|\,\mathit{claimHash}\,\|\,k\,\|\,\mathrm{MTroot}(\Ind(i))\,\|\,\vk_i\,\|\,x)$
(rule S-4 of \cref{subsec:sec-fs}). This answers (i) and (iii). It does not answer (ii), which
breaks the strong transform; that attack is part of the named deployment assumption of
\cref{subsec:rec-cannot}.

\paragraph{Sampler bias.} \Cref{thm:bg-bcs} treats challenges as uniform elements of $C_j$. An
implementation maps hash bits to field elements, and any fixed bad set gains at most the
statistical distance of that map from uniform, which therefore adds directly to each $\eps_j$.
Over a prime field, reducing a $b$-bit uniform integer modulo $p$ (write $2^b=Qp+r$, $0\le r<p$)
gives distance $r(p-r)/(p\,2^b)\le p/2^{b+2}$. For a 31-bit prime and $b=128$ this is about
$2^{-99}$, which would dominate a $2^{-128}$ round error, so implementations need rejection
sampling or $b\ge160$ bits per coordinate. Over $\mathrm{GF}(2^n)$ every $n$-bit string is an
element, so sampling is exactly uniform (\cref{subsec:bg-binary}).

\begin{takeaway}[What is and is not provable today]
The information-theoretic statement (completeness and RBR-KS of a holographic IOP) is the only
link that submissions can prove in Lean with current libraries. ArkLib has no BCS transform and
no Fiat--Shamir soundness theorem in its legacy layer; its Fiat--Shamir completeness lemma is
\lean{sorry}. The compiled bound is therefore a cited meta-theorem with the explicit form
$(t+k)\,\eps_{\mathrm{rbr}}+4.5\,\max(t,t^*)^2/2^\lambda$, with $t^*$ the side condition's threshold
(\cref{subsec:sec-cap}). Weak-FS and generated-statement attacks are excluded by contract rules,
not by theorems; the hash-in-circuit attack is excluded by neither and is a named deployment
assumption.
\end{takeaway}

%-----------------------------------------------------------------------------------------------
\subsection{Proximity gaps: what is proven and what has been refuted}\label{subsec:bg-proximity}

Let $C=\mathrm{RS}[\F,D,k]$ be the Reed--Solomon code of evaluations on $D\subseteq\F$, $|D|=n$, of
polynomials of degree $<k$. Its rate is $\rho=k/n$ and its relative distance is
$\delta_C=1-\rho+1/n$. Write $\Delta(u,C)$ for the relative Hamming distance from $u\in\F^D$ to
$C$. Three radii matter: \emph{unique decoding} $\gamma<\delta_C/2\approx(1-\rho)/2$, where every
word has at most one codeword within $\gamma$; the \emph{Johnson radius} $1-\sqrt\rho$, up to which
lists are provably polynomial; and \emph{capacity} $1-\rho$.

\begin{definition}[Proximity gap; correlated agreement]\label{def:bg-ca}
The code $C$ has \emph{correlated agreement for lines} at radius $\gamma$ with exceptional count
$a$ if, for all $u_0,u_1\in\F^D$,
\begin{multline*}
  \bigl|\{z\in\F : \Delta(u_0+z\,u_1,\,C)\le\gamma\}\bigr|> a \\
  \Longrightarrow\ \exists S\subseteq D,\ |S|\ge(1-\gamma)n,\ \exists v_0,v_1\in C:\
  u_0|_S=v_0|_S,\ u_1|_S=v_1|_S.
\end{multline*}
The combination step of a batched or folded test then errs with probability about $a/|\F|$.
\emph{Mutual correlated agreement} (MCA), which WHIR needs \cite{ACFY24-WHIR}, additionally
requires the agreement set of the random combination itself to be one on which $u_0,u_1$ agree
with codewords, except with similar probability over $z$.
\end{definition}

The soundness error of FRI, STIR and WHIR decomposes into three terms. The batching or folding term
is proportional to $a/|\F|$ per round. The query term is about $(1-\gamma)^s$ for $s$ queries. The
third term comes from Fiat--Shamir (\cref{subsec:bg-bcs}). The radius $\gamma$ at which the gap
is \emph{proven} therefore sets the number of queries, and query count is the dominant verifier
cost.

\begin{table}[t]
\centering\small
\begin{tabularx}{\linewidth}{@{}l>{\raggedright\arraybackslash}Xl@{}}
\toprule
Regime & Result (exceptional count $a$; error $\approx a/|\F|$) & Status \\
\midrule
$\gamma<\delta_C/2$ (UD) & correlated agreement with $a\approx n$ \cite{BCIKS20} & proven \\
$\gamma\in(\delta_C/3,\delta_C/2-1/n)$ & $a=O(1)$ with arbitrarily small proximity loss \cite[Thm.~1.3]{BCHKS25} & proven \\
$\gamma<1-\sqrt\rho-\eta$ (JB) & $a=O(n^2/\eta^7)$ \cite{BCIKS20}; improved to $a=O(n/\eta^5)$ \cite[Thm.~1.5]{BCHKS25} & proven \\
MCA up to Johnson & for all polynomial generators (lines, powers, curves) \cite{BCGM25-MCA}; see also \cite{Hab25-MCA} & proven \\
past Johnson by $O(1)$ steps & MCA with relative gain $O(1/n)$ \cite{Jo26-MCA} & preprint \\
up to $1-\rho-\eta$ & optimal gaps for folded RS, multiplicity, subspace-design and \emph{random} RS codes \cite{GG25,GGSW26}; not smooth-domain RS & proven \\
\midrule
up to capacity, RS & BCIKS20 capacity conjecture: refuted for rate $\to0$, radius $\to1$ \cite{DG25-Dist} & refuted \\
up to capacity, RS & CA, WHIR MCA and DEEP-FRI list-decoding conjectures refuted in general form \cite{CS25-PGC} & refuted \\
at Johnson ($\eta=0$), char.~2, rate $\approx1/16$ & $n^{2-o(1)}$ exceptional points on random $\F_2$-subspaces \cite[Thm.~1.6, Cor.~1.7]{BCHKS25}; beyond-Johnson refuted for these codes & lower bound \\
DEEP-ALI, ethSTARK toy problem & attacks showing soundness is governed by list decodability \cite{BCHKS25} & refuted \\
near capacity, prime fields & lines with $2^{\Omega(1/\eta)}$ close points $z$ at $1-\rho-\eta$, $\eta=\Theta(1/\log n)$, on multiplicative subgroups \cite{KKH26}; $\ge n^{C}$ close points for primes $p\equiv1\bmod n$ at $\delta=(1-\rho)-\Omega(1/\log n)$ \cite{Kam26} & refuted \\
\bottomrule
\end{tabularx}
\caption{Proximity gaps for Reed--Solomon codes as of October 2026. ``Refuted'' means that
counterexamples exist in the stated regime, not that every restricted variant fails; open problems
are surveyed in \cite{ABF26}.}
\label{tab:bg-proximity}
\end{table}

Three consequences follow (\cref{tab:bg-proximity}). First, the ``up to capacity'' conjectures on
which most deployed parameter sets rest are false in at least some regimes. Crites and Stewart
propose a weakened ``list-decoding capacity'' conjecture, for prime $q$ up to the radius
$\delta$ with $H_q(\delta)\le1-\rho-\eta$, whose loss against capacity is at most
$1/\log_2q+1/n$ \cite{CS25-PGC}. Second, BCHKS25 show that improved proximity
gaps imply improved list decoding for all RS codes, so beyond-Johnson gaps first require
list-decoding progress \cite{BCHKS25}. Third, the capacity-achieving results of \cite{GG25,GGSW26}
cover codes that FRI-style protocols do not currently use, which makes the choice of code a
legitimate design axis. The Ethereum Foundation's targets count only unique-decoding and Johnson
parameters as provable \cite{EF25-zkEVMsec,soundcalc}. The EF Proximity Prize poses its grand MCA
and list-decoding challenges at $\eps^\ast=2^{-128}$ for rates $1/2$ to $1/16$ \cite{ProximityPrize}.

\paragraph{Cost in queries.} Each query buys $-\log_2(1-\gamma)$ bits. At $\lambda=128$, ignoring
the slack $\eta$, list-size factors and grinding, this gives \cref{tab:bg-queries}. The Johnson
column is therefore a lower bound.

\begin{table}[t]
\centering\small
\begin{tabular}{@{}lccccc@{}}
\toprule
& \multicolumn{2}{c}{UD, $\gamma=(1-\rho)/2$ (proven)} & \multicolumn{2}{c}{JB, $\gamma\approx1-\sqrt\rho$ (proven)} & CB, $\gamma\approx1-\rho$ (refuted) \\
\cmidrule(lr){2-3}\cmidrule(lr){4-5}\cmidrule(l){6-6}
$\rho$ & bits/query & queries & bits/query & queries & queries \\
\midrule
$1/2$  & 0.415 & 309 & 0.5 & 256 & 128 \\
$1/4$  & 0.678 & 189 & 1.0 & 128 & 64 \\
$1/8$  & 0.830 & 155 & 1.5 & 86  & 43 \\
$1/16$ & 0.913 & 141 & 2.0 & 64  & 32 \\
\bottomrule
\end{tabular}
\caption{Queries $s=\lceil 128/(-\log_2(1-\gamma))\rceil$ needed for a $2^{-128}$ query term,
derived from the radii (R7 reports $\approx257$ for JB at $\rho=1/2$ once slack is included).
Proven Johnson parameters cost exactly $2\times$ the capacity-conjectured queries (bits per query
halve), and unique decoding $2.4$--$4.4\times$.}
\label{tab:bg-queries}
\end{table}

\paragraph{Status in Lean.} ArkLib proves the BCIKS20 unique-decoding correlated agreement for
affine lines without \lean{sorry}. Its Johnson-regime BCIKS20 theorem (Thm.~5.1) is \lean{sorry}.
A separate RS development proves a Johnson-range MCA bound for lines without
\lean{sorry}.\footnote{\repo{ArkLib@db7005f}: \path{ArkLib/Data/CodingTheory/ProximityGap/BCIKS20/AffineLines/UniqueDecoding.lean:33} (proved);
\path{.../BCIKS20/AffineLines/Main.lean:38} (\lean{sorry}); \lean{mcaError\_affineLine\_johnson\_le},
\path{ArkLib/Data/CodingTheory/ReedSolomon/MutualCorrelatedAgreement/Johnson/Probability.lean:41} (proved).} The BCHKS25 Johnson-range MCA
bound and some list-decoding bounds are admitted with \lean{sorry}. A Lean-checked competition
with an axiom whitelist therefore forces proven regimes, either proved by the submitter or
supplied as a named, human-audited hypothesis transcribed from a published theorem. The Johnson
regime also enlarges the batching term $a/|\F|$ with $a=O(n/\eta^5)$, which, with one challenge
per check, forces large challenge fields: extensions of degree $4$--$6$ over a 31-bit prime field,
and $\mathrm{GF}(2^{192})$ rather than $\mathrm{GF}(2^{128})$ in the binary case
(\cref{subsec:bg-binary}).

%-----------------------------------------------------------------------------------------------
\subsection{Binary fields for SNARKs}\label{subsec:bg-binary}

recursion.fast fixes its relation over a binary field (see the protocol-model section). This
subsection collects what that choice implies for codes, costs, soundness and Lean support. The
numbers in \cref{tab:bg-binary-errors} and the cycle counts come from the design's own arithmetic
and probes, which apply the cited theorems' formulas; they are derived figures, not results
quoted from the literature.

\paragraph{Fields and towers.} $\mathrm{GF}(2^n)=\mathrm{GF}(2)[x]/(f)$ for an irreducible $f$ of
degree $n$. An element is a polynomial of degree $<n$, stored as an $n$-bit word with bit $i$
holding the coefficient of $x^i$. Addition is XOR, so $a+a=0$ and subtraction equals addition.
Multiplication is a \emph{carry-less} product of degree $<2n$ followed by reduction modulo $f$.
Binius \cite{DP23-Binius} uses Wiedemann-style towers $\mathcal T_{i+1}=\mathcal T_i[X]/(X^2+X_iX+1)$,
in which every $\mathcal T_i$ is a subfield of $\mathcal T_{i+1}$ and subfield products are cheap.
recursion.fast uses leanVM's pair \cite{leanVM}:
\[
  K=\mathrm{GF}(2)[x]/(x^{64}+x^4+x^3+x+1)=\mathrm{GF}(2^{64}),\qquad
  E=K[y]/(y^3+y+1)=\mathrm{GF}(2^{192}).
\]
The cubic $y^3+y+1$ is irreducible over $K$. It is irreducible over $\mathrm{GF}(2)$, so a root
would generate $\mathrm{GF}(8)$ and have multiplicative order $7$, but $7\nmid 2^{64}-1$ because
$2^{64}\equiv2\pmod 7$. Every 64-bit word is an element of $K$ and every 24-byte string is an
element of $E$. Encodings are therefore canonical, and Fiat--Shamir sampling from 24 uniform hash
bytes is \emph{exactly} uniform on $E$, so the sampler-bias issue of \cref{subsec:bg-bcs} does not
arise.

\paragraph{Additive FFTs: no 2-adicity constraint.} The multiplicative group of
$\mathrm{GF}(2^n)$ has odd order $2^n-1$, so it has no subgroups of order $2^k$ and the
multiplicative FFT domains of FRI, STIR and WHIR do not exist. Reed--Solomon codes are instead
evaluated on $\F_2$-linear subspaces (or their affine cosets) $S\subseteq K$ of size $2^{d}$,
$d\le 64$. The Lin--Chung--Han additive FFT evaluates a polynomial of degree $<h$ on such an $S$
in $(h/2)\log_2h$ multiplications and $h\log_2h$ additions, using a ``novel polynomial basis''
\cite{LCH14-AdditiveFFT}. Every power-of-two domain size up to $2^{64}$ therefore exists at every
rate. Over a prime field, by contrast, the largest smooth domain is $2^{v_2(p-1)}$: $2^{27}$ for
BabyBear and $2^{24}$ for KoalaBear. Mapping a query index to a domain point is also free in the
binary case (the index bits select basis vectors), where a prime field pays an exponentiation.
Folding over characteristic~2 uses the line $u_0+r(u_0+u_1)$, an invertible change of
coordinates from the affine line of \cref{def:bg-ca}.

\paragraph{Carry-less multiplication and Zbc.} The basic operation is the $64\times64\to128$-bit
carry-less product. RISC-V provides it in the Zbc extension, with \texttt{clmul} (low half),
\texttt{clmulh} (high half) and \texttt{clmulr} (reversed) \cite{RISCV-Bitmanip}. The design's
probes measure a $K$ multiplication on RV64IM at 379 cycles with 4-bit tables and no Zbc, against
35 cycles with Zbc. An $E$ multiplication (three-term Karatsuba, six carry-less pairs) costs 103
instructions, or 139 weighted units with \texttt{clmul} at weight 4. This is why the competition's machine is RV64IM+Zbc (see the
architecture and cost-model sections).

\paragraph{Binary-field proof systems.}
\begin{itemize}[nosep]
\item \emph{Binius} \cite{DP23-Binius} commits to multilinears over towers with a Brakedown-style
  code.
\item \emph{FRI-Binius} \cite{DP24-FRIBinius} replaces that commitment with FRI over additive
  domains. It also introduces \emph{ring switching} with \emph{packing}: $2^\kappa$ elements of a
  small field (for example bits) are packed into one element of a degree-$2^\kappa$ extension, and
  a claim about the small-field multilinear is reduced to one claim about the packed multilinear.
  Committing $2^\kappa$ bits then costs about as much as committing one large-field element.
\item \emph{Binius64} \cite{Binius64} is the current Binius design, built around 64-bit words.
\item \emph{Ligerito} \cite{NA25-Ligerito} is a Ligero-style commitment repeated with sumcheck.
  It is benchmarked over $\F_{2^{128}}$ with $\F_{2^{32}}$ as the first-round field.
\item \emph{Blaze} \cite{BCFRRZ25-Blaze} uses repeat--accumulate--accumulate (RAA) codes over binary
  fields with code switching into BaseFold. It has a fast prover, but needs about 1{,}059 column
  queries at 100 bits and a verifier of about 150\,ms.
\item \emph{Flock} \cite{BRW26-Flock} proves large batches of Boolean computations (SHA-256, Keccak,
  BLAKE3), and leanVM uses it to prove its BLAKE2s calls \cite{leanVM}.
\end{itemize}
The published binary systems we checked target about 100 bits over 128-bit fields: FRI-Binius's
provable analysis is at 100 bits over $\F_{2^{128}}$, and Ligerito is at $\lambda=100$
\cite{DP24-FRIBinius,NA25-Ligerito}. leanVM's current release instead claims 128 bits over
$\mathrm{GF}(2^{192})$ in the Johnson regime, with no proximity-gap conjecture \cite{leanVM}.

\paragraph{Why $\mathrm{GF}(2^{192})$.} With a \emph{single} challenge per round, every round
error has the form (numerator)$/|E|$ with numerator at least $2$: the degree $d\ge2$ of a sumcheck
round, the $(J-1)L$ of a batching step, the $\gamma n+1$ of a unique-decoding proximity step
\cite[Cor.~1.4]{BCHKS25}, or the exceptional count $a$ of a Johnson-regime MCA step. So one
$\mathrm{GF}(2^{128})$ element cannot meet a $2^{-128}$ cap. In-round repetition (several
independent challenges in one \lean{chal} round, as Plonky2 does over Goldilocks) would meet it at
roughly the same field cost, since two $\mathrm{GF}(2^{128})$ products cost as many $K$-products as
one $\mathrm{GF}(2^{192})$ product, but it doubles the affected prover work and oracle answers and
complicates every protocol. The design therefore fixes $E=\mathrm{GF}(2^{192})$, with one challenge
per check, as the simpler option, not the only one. \Cref{tab:bg-binary-errors} shows the margins
for a $2^{24}$-element witness. At $e=192$ the Johnson regime is reachable with slack
$\eta\approx0.011$--$0.018$, i.e.\ 268--276 queries at $\rho=1/2$ instead of the slack-free 256 of
\cref{tab:bg-queries} ($+5$--$8\%$). A 256-bit field would recover this slack at about $1.5\times$
the multiplication cost.

\begin{table}[t]
\centering\small
\begin{tabularx}{\linewidth}{@{}>{\raggedright\arraybackslash}X>{\raggedright\arraybackslash}p{3.6cm}ccc@{}}
\toprule
Round type & Error & $|E|=2^{128}$ & $|E|=2^{192}$ & $|E|=2^{256}$ \\
\midrule
sumcheck, degree $d\le3$ & $d/|E|$ & $2^{-126.4}$ & $2^{-190.4}$ & $2^{-254.4}$ \\
batching $J\le2^9$ claims, list size $L\le2^6$ & $(J-1)L/|E|$ & $2^{-113}$ & $2^{-177}$ & $2^{-241}$ \\
UD proximity, $\gamma\to(1-\rho)/2$, $n=2^{25}$ \cite[Cor.~1.4]{BCHKS25} & $(\gamma n+1)/|E|$ & $2^{-105}$ & $2^{-169}$ & $2^{-233}$ \\
Johnson MCA, $\rho=1/2$, $\eta=0.0177$ (276 queries) \cite[Thm.~4.6]{BCHKS25} & $2^5a/|E|$, $a\approx2^{52.6}$ & fails & $2^{-134}$ & $2^{-198}$ \\
Johnson MCA, $\rho=1/2$, $\eta=0.011$ (268 queries) & $2^5a/|E|$, $a\approx2^{56}$ & fails & $2^{-131}$ & $2^{-195}$ \\
Johnson MCA, $\rho=1/2$, $\eta=0.0055$ (262 queries) & $2^5a/|E|$, $a\approx2^{61}$ & fails & $2^{-126}$ & $2^{-190}$ \\
Johnson MCA, $\rho=1/4$, $\eta=0.0125$ (133 queries) & $2^5a/|E|$, $a\approx2^{55}$ & fails & $2^{-132}$ & n/c \\
\bottomrule
\end{tabularx}
\caption{Per-round errors against a $2^{-128}$ cap for a $2^{24}$-element witness, one challenge
per round, at block length $n=2^{24}/\rho$, derived in the design (D3 §1.2) from the cited bounds.
The UD row is the limit $\gamma\to(1-\rho)/2$ of BCHKS25 Cor.~1.4, which holds for
$\gamma<\delta/2-3/(\delta n)$. $a$ follows BCHKS25 Thm.~4.6 as transcribed by leanVM, with the
conservative $m=\max(\lceil\sqrt\rho/\eta\rceil,3)$ (Thm.~1.5's $\max(\lceil\sqrt\rho/(2\eta)\rceil,3)$
would be about 5 bits smaller), multiplied by leanVM's folding factor $2^{\ell-1}=2^5$ for fold
factor $\ell=6$; this factor is not part of Thm.~4.6 for one line (its $M$-function form multiplies
by $M$), and we have not re-derived the generator it corresponds to. n/c: not computed. A single
$2^{128}$ challenge misses the cap in every row, and $2^{192}$ meets it in the Johnson regime when
$\eta\gtrsim0.011$.}
\label{tab:bg-binary-errors}
\end{table}

\paragraph{Proximity gaps over binary fields.} The positive results of
\cref{tab:bg-proximity} are field-agnostic. BCIKS20's correlated-agreement theorems are proved
for a general field and any evaluation set $D\subseteq\F$ \cite{BCIKS20}. The Johnson-regime
bounds of BCHKS25 (Thms.~1.5 and 4.6) are stated with no characteristic restriction, but the MCA
version (Thm.~4.6) has only a proof sketch for the separable case; the inseparable case, which is
the relevant one in characteristic~2, is argued only for the line version (Thm.~1.5), by reference
to BCIKS20 App.~C \cite{BCHKS25}. BCGM25's MCA up to Johnson, a refereed result, holds over any
finite field \cite{BCGM25-MCA}. These bounds therefore apply to RS codes on $\F_2$-subspaces, with
the caveat on Thm.~4.6 recorded in \cref{subsec:sec-hypotheses}. The characteristic-2 negative
results are BCHKS25 Thm.~1.6 and Cor.~1.7. For every field of characteristic~2 there are RS codes
on (essentially random) $\F_2$-subspaces, of rate about $1/16$, with $n^{2-o(1)}$ exceptional points
at the Johnson radius ($\eta=0$) and a constant proximity loss; at higher rates the same family
refutes the capacity conjecture below capacity. So no bound linear in $n$ can hold at $\eta=0$ for
all such codes, and a beyond-Johnson gap over characteristic~2 would have to depend on the choice of
domain (the paper leaves this open). Nothing is implied for $\eta>0$, which is exactly the
proven-regime policy, and \cref{tab:bg-binary-errors} already prices the slack; whether the
$1/\eta^5$ dependence is tight is open.

\paragraph{Lean support.} CompPoly defines $K$ as \lean{BF64}, a wrapper around
\lean{BitVec 64}. Multiplication is reduction composed with a shift--XOR carry-less product, and
it is proved to agree with Mathlib's \lean{AdjoinRoot} model of $\mathrm{GF}(2)[x]/(f)$.
Irreducibility of the degree-64 modulus is proved by Rabin's test: $X^{2^{64}}\equiv X \bmod f$
and a coprimality check at $2^{32}$. A generated 552-line certificate is checked by \lean{rfl}
in the kernel, with no \lean{native\_decide} and no \lean{sorry}; the file compiles in about 88\,s.
The cubic extension's irreducibility follows from the order argument above.\footnote{\repo{CompPoly@572f997},
\path{CompPoly/Fields/Binary/BF64/Basic.lean:130} (\lean{basePoly\_irreducible}) and
\path{CompPoly/Fields/Binary/BF64/Ext3.lean:127} (\lean{ext3Poly\_irreducible}).} Unpatched
comparator replays a challenge that depends on both certificates in 101\,s. This contrasts with the
Proximity Prize, whose KoalaBear sextic certificate forced a comparator patch (\cref{sec:prior}).
CompPoly also proves the LCH14 additive NTT correct, though concrete instances exist only for
towers. Mathlib's \lean{GaloisField 2 64} is unusable for the relation: it is a splitting field
with no basis, no decidable equality and no computable arithmetic. ArkLib's gaps over binary fields
are substantial:
\begin{itemize}[nosep]
\item Its Reed--Solomon definitions accept any domain, and its unique-decoding correlated agreement
  and Johnson-range line MCA have no characteristic hypothesis and no \lean{sorry}.
\item Its FRI requires \lean{[NonBinaryField F]}, and WHIR's constrained RS codes need
  multiplicative smooth domains. There is no formal FRI or WHIR over binary
  fields.\footnote{\repo{ArkLib@db7005f}, \path{ArkLib/ProofSystem/Fri/Spec/General.lean:37}.}
\item Its capacity-regime MCA needs characteristic $0$ or above $k-1$, so it fails in characteristic
  2 for $k>2$.
\item Its Binius, FRI-Binius and ring-switching developments are written over abstract
  characteristic-2 fields, and every completeness and RBR-KS theorem in them depends on
  \lean{sorry} (23, 9 and 14 sorries in BinaryBasefold, FRIBinius and RingSwitching packing).
\item Its native sumcheck soundness holds over any finite field.
\end{itemize}
The cheapest proven path for a binary-field submission today is therefore a unique-decoding RS (or
Ligero-style) IOP over $E$ on an additive domain inside $K$, combined with sumcheck. This is our
inference from the available lemmas, not a demonstrated development.

%-----------------------------------------------------------------------------------------------
\subsection{Hash-based SNARK verifiers and their costs}\label{subsec:bg-costs}

Hash-based verifiers spend their time on three things: hash compressions for Merkle paths
(proportional to queries $\times$ depth), field arithmetic in an extension field (sumcheck rounds,
folding, out-of-domain checks), and moving the proof's bytes. \emph{FRI}
\cite{BBHR18-FRI,BGKS20-DEEPFRI} tests proximity to RS codes with $O(\lambda\log d)$ queries.
\emph{STIR} \cite{ACFY24-STIR} shifts the code rate across rounds and needs
$O(\log d+\lambda\log\log d)$ queries. \emph{WHIR} \cite{ACFY24-WHIR} gives a multilinear
commitment with RBR-KS whose verifier performs no field divisions. \emph{BaseFold}
\cite{ZCF24-BaseFold} uses foldable codes and is proven in the unique-decoding regime.
\emph{FRI-Binius} \cite{DP24-FRIBinius} works over binary towers, superseding the PCS of
\cite{DP23-Binius}. \emph{Ligerito} \cite{NA25-Ligerito} is a Ligero-style \cite{AHIV17-Ligero}
scheme repeated with sumcheck. Brakedown \cite{GLSTW23-Brakedown} and Orion \cite{XZS22-Orion}
have linear-time provers and verifiers costing $O(\sqrt N)$, which makes them poor fits for a cheap
verifier. \emph{Circle STARKs} \cite{HLP24-Circle} run FRI over the circle curve of the Mersenne
prime $2^{31}-1$. Binary-field systems (Binius, FRI-Binius, Ligerito, Blaze) are described in
\cref{subsec:bg-binary}.

\begin{table}[t]
\centering\footnotesize
\begin{tabularx}{\linewidth}{@{}>{\raggedright\arraybackslash}p{4.2cm}>{\raggedright\arraybackslash}p{1.55cm}>{\raggedright\arraybackslash}X>{\raggedright\arraybackslash}p{2.0cm}@{}}
\toprule
System, setting & Field & Verifier cost & Regime \\
\midrule
WHIR-JB, $2^{24}$ vars, $\rho=1/2$ & prime: Goldilocks$^3$ & 299\,KiB, $\approx5.1$k hashes, 2.5\,ms \cite[Tab.~5]{ACFY24-WHIR} & proven (JB, via \cite{BCGM25-MCA}); parameters as published, not re-derived under BCGM25's constants \\
WHIR-UD, same & prime: Goldilocks$^3$ & 621\,KiB, $\approx10$k hashes, 4.8\,ms & proven (UD) \\
WHIR-CB, same & prime: Goldilocks$^3$ & 156\,KiB, $\approx2.7$k hashes, 1.4\,ms & conjectured \\
WHIR at $\rho=1/16$, $2^{24}$ (UD/JB/CB) & prime: Goldilocks$^3$ & 535/140/74\,KiB; 10k/2.7k/1.4k hashes & as above \\
FRI/STIR/WHIR, degree $2^{24}$, $\rho=1/2$ & prime, 192-bit & 306/160/157\,KiB; 5.6k/2.6k/2.7k hashes; 3.9/3.8/1.0\,ms \cite[Tab.~3]{ACFY24-WHIR} & conjectured \\
ethSTARK, $\rho=1/4$, 20-bit grinding, 128 bits & prime & 140 queries (provable) vs 55 (conjectured) \cite{ethSTARK21} & both \\
FRI-Binius, $2^{24}$ $\F_2$-vars, $\rho=1/2$, 100 bits & binary, $\F_{2^{128}}$ & 0.202\,MiB \cite{DP24-FRIBinius} & proven \\
Ligerito, $2^{24}$, $\rho=1/4$, 100 bits & binary, $\F_{2^{128}}$ & 255\,KiB, 148 queries per round \cite{NA25-Ligerito} & UD (derived) \\
\midrule
Plonky2 recursion circuit & prime: Goldilocks & $2^{12}$ gates, $\approx75\%$ Merkle hashing \cite{Plonky2} & 100 bits, conjectured \\
leanVM 2$\to$1 recursion, 128 bits & binary, $\mathrm{GF}(2^{192})$ & $562{,}737$ leanISA cycles, 190\,KiB, 0.27\,s (live README); $686{,}085$ cycles, 206\,KiB, 0.39\,s at commit \texttt{a386121} \cite{leanVM} & proven (JB, no gap conjecture) \\
leanVM (KoalaBear branch), recursion & prime: KoalaBear$^5$ & proven vs conjectured: $\approx1.4\times$ time, $1.9$--$2\times$ size \cite{leanVM} & both \\
Jolt verifier in a Jolt RISC-V guest & prime (128-bit) & $6.8$--$7.7\times10^{7}$ cycles, with field and hash inline instructions (a16z/jolt draft PRs \#1892, \#1911, Sep.~2026) & lattice PCS \\
Airbender recursion verifier (RISC-V) & prime: M31$^4$ & run with a $10^7$-cycle budget; an upper bound only & 64--67 bits provable \cite{soundcalc} \\
SP1 Hypercube recursion & prime: KoalaBear$^4$ & 124 queries at $\rho=1/4$, 16-bit PoW \cite{soundcalc} & 100 bits UD (proven) \\
\bottomrule
\end{tabularx}
\caption{Concrete verifier-cost anchors, labelled by field: ``prime'' rows are prime-field
systems, kept for comparison, and do not transfer directly to recursion.fast's binary fields.
Proof sizes and hash counts are the papers' or repositories' figures; ``regime'' records whether
the soundness parameters rest on a proven proximity bound. Rows below the rule are recursion
steps.}
\label{tab:bg-costs}
\end{table}

\Cref{tab:bg-costs} collects anchors near 128 bits. Five observations matter for the competition.
(1) The WHIR-JB column is now the best \emph{proven} WHIR data point, because the MCA-to-Johnson
conjecture it assumed is a theorem \cite{BCGM25-MCA}, with the caveat that BCGM25's exceptional
counts are quadratic in $n$ and the published parameters have not been re-derived under them. The CB column has no proof. (2) The cleanest
public measurement of the price of provability on a real recursion step is leanVM's KoalaBear
branch: about $1.4\times$ the time and $1.9$--$2\times$ the size. (3) leanVM's 2-to-1 step costs
about $5.6\times10^5$ cycles of a field-native ISA over $\mathrm{GF}(2^{192})$ with a hash
instruction, about $120$--$140\times$ fewer
than the Jolt in-guest figure. (4) We found no published cycle count for a STARK verifier running
as \emph{plain} RV64IM guest code. SP1 and RISC Zero defer in-guest verification to dedicated
recursion circuits. We estimate $10^7$--$10^8$ cycles without precompiles (an estimate, not a
measurement). (5) Precompiles (hash, field multiplication, carry-less multiplication) change
rankings by $10$--$100\times$ (Jolt inlines, Airbender delegation, leanVM's native ISA), so the
instruction-set and precompile policy is a first-order design decision (\cref{sec:cost}).
Production parameter sets vary widely in provable security. Soundcalc reports, for example,
128-bit Johnson security for ZisK and 100-bit unique-decoding security for SP1 Hypercube, while
RISC Zero's documented parameters rest on the toy-problem conjecture \cite{soundcalc}.

%-----------------------------------------------------------------------------------------------
\subsection{Recursion, aggregation and accumulation}\label{subsec:bg-recursion}

\paragraph{IVC and PCD.} \emph{Incrementally verifiable computation} (IVC) \cite{Val08} proves
$z_n=F^{(n)}(z_0)$ step by step, with a proof whose size is independent of $n$. \emph{Proof-carrying
data} (PCD) \cite{CT10,BCCT13} generalises IVC from paths to directed acyclic graphs: every vertex
checks a \emph{compliance predicate} $\varphi$ on its output, local data and incoming messages.
Knowledge soundness asks for an extractor that outputs a $\varphi$-compliant transcript whenever
the final proof verifies (formalised in \cite[§5]{BCMS20}). There are two ways to build PCD.
\emph{Full recursion} verifies the child proofs inside the next proof, so the recursion circuit
contains the whole SNARK verifier. \emph{Accumulation} (or folding) defers the expensive part.

\begin{definition}[Split accumulation {\cite{BCLMS21}}; folding {\cite{KST22-Nova}}]\label{def:bg-acc}
Let $\Phi$ be a predicate on inputs $q=(q.x,q.w)$, split into instance and witness parts. A
\emph{split accumulation scheme} consists of a prover $\mathsf{P}$, a verifier $\mathsf{V}$ and a
decider $\mathsf{D}$. Accumulators also split, as $\mathrm{acc}=(\mathrm{acc}.x,\mathrm{acc}.w)$.
The prover maps $(q_1,\dots,q_n)$ and old accumulators $(\mathrm{acc}_1,\dots,\mathrm{acc}_m)$ to a
new accumulator $\mathrm{acc}$ and a proof $\pi$. The verifier sees \emph{instance parts only}:
$\mathsf{V}([q_i.x],[\mathrm{acc}_j.x],\mathrm{acc}.x,\pi)\in\{0,1\}$. The scheme must satisfy two
properties.
\begin{itemize}[nosep]
\item \emph{Completeness:} if $\Phi(q_i)$ holds for all $i$ and $\mathsf{D}(\mathrm{acc}_j)$ for all
  $j$, then $\mathsf{V}$ accepts and $\mathsf{D}(\mathrm{acc})$ holds.
\item \emph{Knowledge soundness:} if $\mathsf{V}$ accepts and $\mathsf{D}(\mathrm{acc})$ holds, an
  extractor outputs $[q_i.w]$ and $[\mathrm{acc}_j.w]$ such that all $\Phi(q_i)$ and all
  $\mathsf{D}(\mathrm{acc}_j)$ hold.
\end{itemize}
A \emph{folding scheme} is the special case $\mathsf{D}=\Phi$, i.e.\ an interactive protocol mapping
two instance--witness pairs of $\R$ to one pair of $\R$ \cite{KST22-Nova}.
\end{definition}

The defining invariant is that a $\mathsf{D}$-valid accumulator certifies everything absorbed so
far, so the recursion circuit runs only $\mathsf{V}$ and never $\mathsf{D}$. Batching and
aggregation, in contrast, compress a \emph{fixed} set of finished proofs. Accumulation and
folding are both instances of a \emph{reduction of knowledge} \cite{KP23-RoK}.

\begin{definition}[Reduction of knowledge, information-theoretic]\label{def:bg-rok}
A reduction from $\R_{\mathrm{in}}$ to $\R_{\mathrm{out}}$ is a public-coin interactive protocol in
which the prover holds $(s,w)$ and the verifier holds $s$. At the end the verifier outputs
$s'$ (or rejects) and the prover outputs $w'$. It is \emph{perfectly complete} if
$(s,w)\in\R_{\mathrm{in}}$ implies, for every challenge vector, that the verifier outputs some
$s'$ with $(s',w')\in\R_{\mathrm{out}}$. It is \emph{RBR knowledge sound} with errors $(\eps_m)$ if
\cref{def:bg-rbrks-warp} holds with $\R=\R_{\mathrm{in}}$ in the (empty) law and
$\R_{\mathrm{out}}$ in the (full) law. A proof system is the case
$\R_{\mathrm{out}}=\{(\top,\cdot)\}$.
\end{definition}

An accumulation scheme for R1CS in this model is three reductions, which together form the
recursion tree of \cref{fig:bg-tree}. $\mathsf{embed}$ reduces $\Rrcs$ to $\Racc$; $\mathsf{fold}$
reduces $\Racc\times\Racc$ to $\Racc$; $\mathsf{decide}$ reduces $\Racc$ to $\top$. RBR errors
compose additively under sequential composition. The accumulator statement must have bounded size,
or ``batching by concatenation'' is trivially sound and complete. WARP's Thm.~2.8 (formally B.4) turns RBR-KS of the reduction into straight-line SR knowledge
soundness with error $(t+k)\cdot\eps_{\mathrm{rbr}}$, and its Thm.~10.3 compiles the result into a
non-interactive accumulation scheme \cite{BCFW25-WARP}.

\begin{figure}[t]
\centering
\begin{tikzpicture}[
  every node/.style={font=\small},
  leaf/.style={draw,rounded corners=2pt,fill=black!3,inner sep=2pt},
  acc/.style={draw,circle,inner sep=1.2pt,fill=blue!6},
  op/.style={font=\scriptsize\itshape,text=blue!50!black},
  arr/.style={-{Latex[length=1.6mm]}}]
\foreach \j in {1,...,4} { \node[leaf] (r\j) at (2.4*\j,0) {$(i,x_{\j},w_{\j})$}; }
\node[font=\scriptsize,anchor=west] at (10.3,0) {$\in\Rrcs$};
\foreach \j in {1,...,4} { \node[acc] (a\j) at (2.4*\j,1.25) {$\mathrm{acc}_{\j}$}; \draw[arr] (r\j) -- node[op,right]{embed} (a\j); }
\node[acc] (b1) at (3.6,2.6) {$\mathrm{acc}_{12}$};
\node[acc] (b2) at (8.4,2.6) {$\mathrm{acc}_{34}$};
\draw[arr] (a1) -- node[op,above left]{fold} (b1); \draw[arr] (a2) -- (b1);
\draw[arr] (a3) -- node[op,above left]{fold} (b2); \draw[arr] (a4) -- (b2);
\node[acc] (c) at (6,3.9) {$\mathrm{acc}_{1..4}$};
\draw[arr] (b1) -- node[op,above left]{fold} (c); \draw[arr] (b2) -- (c);
\node[draw,rounded corners=2pt,fill=green!8] (d) at (6,5.0) {$\mathsf{decide}$: accept / reject};
\draw[arr] (c) -- (d);
\node[align=left,font=\scriptsize,anchor=west] at (9.6,3.6) {depth $\log_2 N$;\\ each fold verifier\\ sees instance parts only};
\end{tikzpicture}
\caption{A 2-to-1 accumulation tree for $N=4$ R1CS instances. Each edge is a reduction of
knowledge; the cost of recursion is the fold verifier's, and the decider runs once.}
\label{fig:bg-tree}
\end{figure}

\paragraph{Accumulation without homomorphism.} Group-based folding (Nova and successors
\cite{KST22-Nova,KS24-HyperNova,BC23-ProtoStar,EG23-ProtoGalaxy,NDCTB24-Mangrove}) relies on
additively homomorphic commitments and is not post-quantum. B\"unz, Mishra, Nguyen and Wang give
accumulation from any linear code and Merkle commitments in the ROM \cite{BMNW24}. A
spot-checked linearity test only guarantees that the new accumulator is $\delta$-close to the
claimed combination, so each extraction step loses $\delta$. Soundness is therefore proven only up
to a fixed depth $d$, with a decider family $\mathsf{D}_0\subseteq\dots\subseteq\mathsf{D}_d$, and
there is a concrete attack once $d\delta$ exceeds the unique-decoding radius. \emph{Arc}
\cite{BMNW25-Arc} removes the depth bound for RS codes, with $2\lambda/\log(1/\rho)$ Merkle
openings per input and accumulation up to the list-decoding radius $1-\sqrt\rho$. \emph{WARP}
\cite{BCFW25-WARP} achieves linear-time accumulation for any linear code with MCA, extracting by
erasure correction with the per-round backward extractor of \cref{def:bg-rbrks-warp}. \emph{FICS
and FACS} \cite{BMMS25-FICSFACS} accumulate via code switching with linear-time codes in the
unique-decoding regime; WARP notes that FACS lacks a straight-line extractor for codes not known
to have an efficient error-tolerant decoder. Lattice-based folding
\cite{BC24-LatticeFold,BC25-LatticeFoldPlus} is the other post-quantum route. The deciders of Arc
and WARP are linear in the accumulated instance; recursion.fast admits such deciders, bounded only
by the fuel and amortised in the T2 score (\cref{subsec:rec-score}).

\paragraph{Why recursive soundness is not a theorem in the ROM.} A random-oracle SNARK has a
straight-line extractor that reads the prover's query log, and composing straight-line extractors
costs only additive time. Recursion, however, needs a SNARK for the statement ``$\Vrf^H$ accepts'',
a relation that \emph{itself queries the oracle}. Barbara, Chiesa and Guan prove that such
\emph{relativized} succinct arguments do not exist in the ROM, even interactive ones with
computational soundness \cite{BCG24-RelROM}. Practice therefore replaces $H$ by a concrete hash
inside the circuit and \emph{assumes} that the instantiated scheme is extractable. Under that
assumption, standard-model PCD theorems cover only constant depth, because a recursive extractor
of size $|\tilde\Prv|^c$ per layer costs $|\tilde\Prv|^{c^d}$ at depth $d$
\cite[Rmk.~5.3]{BCMS20}. A 2-to-1 tree over $N$ leaves has depth $\log_2N=\omega(1)$. The known
ways around this either change the oracle model or are heuristic:
\begin{itemize}[nosep]
\item \emph{Arithmetized oracles.} Relativized SNARKs exist in signed, low-degree and arithmetized
  random-oracle models, and PCD for arbitrary depth follows in the AROM \cite{CCS22,CCGOS23-AROM}.
  These models are costly to instantiate and unused in practice.
\item \emph{Straight-line PCD bounds.} If a relativized argument has a straight-line extractor
  with error $\kappa_{\mathrm{ARG}}(\lambda,q,s)$ and time $t_{\mathrm{ARG}}$, PCD over $N$ nodes
  has error $\kappa_{\mathrm{ARG}}(\lambda,q,s+O(N\cdot t_{\mathrm{ARG}}))$, with no dependence on
  depth \cite{CGSY24}. For hash-based PCD this rests on a non-succinct idealisation, so it is
  heuristic evidence, not a proof.
\end{itemize}
leanerVM, the Lean formalisation of leanVM, records this as its open problem T6 (``recursive
CTE''). Its inner knowledge-soundness theorems live in the ROM, but its outer recursion-topology
argument composes plain-model extractors, and no justified interface connects the two. It forbids
a global bridging axiom and keeps the end-to-end theorem as a target schema.\footnote{\repo{leanerVM@8aaf8a0},
\path{README.md:35-38}; \path{docs/architecture.md:297-350}.}

\begin{takeaway}[Consequence for the batcher]
What can be proved today is the \emph{one-step} property: completeness and RBR-KS of the
interactive embed, fold and decide reductions in the information-theoretic model, with an
explicit error. By composition, a tree of fixed shape then has an explicit error sum. Soundness of
the deployed recursive SNARK is a named deployment assumption. It is never a hypothesis of a
competition theorem, and the tree depth must be declared and enforced.
\end{takeaway}

%-----------------------------------------------------------------------------------------------
\subsection{Lean~4, VCVio and ArkLib for cryptographers}\label{subsec:bg-lean}

\paragraph{What a kernel-checked proof certifies.} Lean~4 \cite{MU21-Lean4} is a dependently typed
functional language whose propositions are types and whose proofs are terms. Elaboration (tactics,
automation, macros) produces a term, and a small \emph{kernel} type-checks it. A theorem is
accepted iff its proof term type-checks against its statement in an environment of prior
declarations. A kernel-checked proof therefore certifies exactly the \emph{elaborated statement}
relative to the definitions it mentions and the axioms in its dependency closure. It says nothing
about whether those definitions mean what the reader thinks. This is the key point for
cryptographers: \emph{the proof is trustworthy; the statement must be audited}.

\paragraph{The trusted computing base.}
\begin{itemize}[nosep]
\item \emph{Kernel.} Lean's kernel, optionally rechecked by independent kernels (nanoda, lean4lean)
  on an exported environment \cite{lean4export,nanoda,carneiro2024lean4lean}. Kernels are not
  infallible: Lean~4.34.0 fixed three proof-of-\lean{False} routes \cite{lean4340}, and a
  nested-inductive projection bug that yielded an axiom-free $0=1$ was also accepted by nanoda
  \cite{brossard2026leanproj}. Kernel diversity is defence in depth, not proof.
\item \emph{Axioms.} Mathematics in Mathlib uses three: \lean{propext} (propositional
  extensionality), \lean{Quot.sound} (quotients) and \lean{Classical.choice} (global choice, hence
  excluded middle). \lean{\#print axioms t} lists the axioms in the closure of \lean{t}. Any other
  axiom, including a user-declared \lean{axiom}, is a hole.
\item \emph{\lean{sorry}} closes any goal by the axiom \lean{sorryAx}. It is the standard
  placeholder for an unfinished proof and is detected by the axiom audit, transitively: a theorem
  that uses a sorried lemma inherits \lean{sorryAx}.
\item \emph{\lean{native\_decide}} evaluates a decision procedure with the compiler rather than the
  kernel, through the axiom \lean{Lean.ofReduceBool}. It thereby trusts the compiler and every
  \lean{@[implemented\_by]} and \lean{@[extern]} annotation in scope, so competitions refuse it.
  Kernel-checked \lean{decide} is permitted but slow.
\end{itemize}

\paragraph{Definitional versus statement-level attacks.} Attacks on a proof checker
(\lean{sorry}, extra axioms, compiled evaluation, environment manipulation by metaprograms, kernel
bugs) are handled by tooling. The Lean FRO's \emph{comparator} \cite{comparator} builds a trusted
challenge and the candidate in a sandbox, exports both, checks that each target theorem has the
same type and that every constant its statement uses is identical, audits axioms, and replays the
candidate in a fresh kernel. Attacks on the \emph{statement} are not handled by any checker. They
include vacuous hypotheses, a definition left open (a \emph{definition hole}, which comparator
matches only by name and type), an error bound that is trivial at the chosen parameters, and an
existentially chosen object that makes a property easy. Comparator's documentation states that
definition holes must be checked by an additional, possibly human, verifier. The real failures of
Lean-checked competitions and zkVM verification efforts so far were of this second kind
(\cref{sec:prior}).

\paragraph{VCVio: oracle computations.} VCVio \cite{tuma2024vcvio,tuma2026vcvio} models a
probabilistic oracle algorithm as a \emph{free monad} over an oracle specification: a program is a
finite tree whose nodes are queries and whose leaves are results.

\begin{leancode}
-- VCVio: oracle computations and probabilistic programs
def OracleComp (spec : OracleSpec ι) := PFunctor.FreeM spec.toPFunctor
abbrev ProbComp := OracleComp unifSpec          -- only uniform coin queries
-- simulateQ impl oa : interpret every query of `oa` by `impl` (e.g. a lazy random oracle)
-- Pr[ event | oa ] : probability of `event` under the output distribution (in ℝ≥0∞)
\end{leancode}

\noindent An adversary is an arbitrary term of type \lean{OracleComp spec α}. Its local computation
is free and unbounded; only queries carry information, and a query budget is a predicate on the
tree (\lean{IsQueryBound}). A random oracle is a stateful query implementation with a cache, and
\lean{simulateQ} runs a program against it. Probabilities are \lean{ENNReal}-valued. Failure is
added with \lean{OptionT}, and the probability notation counts failure mass as not satisfying the
event, so verifier aborts count as rejections.

\paragraph{ArkLib: protocols and security.} ArkLib \cite{arklib2026} builds on VCVio. A protocol is
described by its message schedule:

\begin{leancode}
structure ProtocolSpec (n : ℕ) where
  dir : Fin n → Direction          -- P_to_V (message) or V_to_P (challenge)
  «Type» : Fin n → Type            -- type of the i-th message
structure Verifier (oSpec) (StmtIn StmtOut : Type) (pSpec : ProtocolSpec n) where
  verify : StmtIn → FullTranscript pSpec → OptionT (OracleComp oSpec) StmtOut
\end{leancode}

\noindent A \lean{Reduction} pairs a prover with such a verifier. An \lean{OracleReduction} lets
the verifier reach prover messages and input oracles only through queries, and a \lean{Proof} is a
reduction whose output statement is \lean{Bool}. The verifier does not sample its own randomness:
the execution draws challenges uniformly, so every protocol is public-coin. Every security
definition is parameterised by the implementation of the shared oracle \lean{oSpec}; the
information-theoretic setting is \lean{oSpec = []ₒ}. The definitions of completeness, soundness,
straight-line knowledge soundness, RBR soundness and RBR-KS (averaged, worst-case, one-shot and
named-extractor forms) contain no \lean{sorry}. Most \emph{general theorems} relating them are
\lean{sorry}, however. This includes RBR-KS $\Rightarrow$ KS, every generic sequential-composition
and context-lifting security theorem, the SR implications and Fiat--Shamir completeness. ArkLib has
no BCS transform and no Merkle trees; VCVio has the latter. 281 declarations are transitively
tainted by \lean{sorryAx} at the examined commit.\footnote{\repo{ArkLib@db7005f},
\path{scripts/axiom_baseline.json}; definitions in
\path{ArkLib/OracleReduction/Security/Basic.lean} and \path{RoundByRound.lean}; sorried theorems in
\path{.../Security/Implications.lean:49,79,193} and
\path{.../Composition/Sequential/Append/Security.lean:47-94}.} Two consequences for a
contract follow. Importing ArkLib definitions is safe, provided an axiom audit rejects any use of
the sorried theorems. Concretely, a toy ``send the witness'' R1CS proof shows knowledge soundness
through \lean{rbrKnowledgeSoundness\_implies\_knowledgeSoundness} and compiles, but
\lean{\#print axioms} reports \lean{sorryAx}. And the contract must fix the oracle model and the
error budget itself, because both are parameters chosen by whoever proves the theorem.

The worst-case RBR-KS of \cref{def:bg-rbrks-warp} reads, in the counting form adopted by the design
(local restatement, with a proved bridge to ArkLib's
\lean{rbrKnowledgeSoundnessWorstCaseWith}):

\begin{leancode}
def KSF.Bad (kSF : KSF relIn relOut V E) (i : pSpec.ChallengeIdx) (s)
    (tr : pSpec.Transcript i.1.castSucc) (c : pSpec.Challenge i) : Prop :=
  ∃ w, ¬ kSF.toFun i.1.castSucc s tr (E.extractMid i.1 s (tr.concat c) w)
       ∧ kSF.toFun i.1.succ s (tr.concat c) w
def RbrKS (kSF : KSF relIn relOut V E) (ε : pSpec.ChallengeIdx → ℝ≥0) : Prop :=
  ∀ s i (tr : pSpec.Transcript i.1.castSucc),
    (Nat.card {c // kSF.Bad i s tr c} : ℝ≥0) ≤ ε i * Fintype.card (pSpec.Challenge i)
\end{leancode}

\noindent The statement involves no \lean{Pr}, \lean{simulateQ} or random-oracle cache: ``at most
an $\eps_i$-fraction of challenges is bad, for every prefix.'' Finiteness of challenge types is
required, because \lean{Nat.card} of an infinite type is $0$, a classic vacuity trap.

\paragraph{What ``information-theoretic knowledge soundness'' means in Lean.} ArkLib's extractors
are Lean functions with no running-time or query bound. ArkLib itself constructs extractors by
\lean{Classical.choose}. Take the extractor that returns a chosen valid witness whenever one
exists. Then straight-line KS collapses to soundness with respect to $L(\R)$, and RBR-KS collapses
to RBR soundness with the same errors. On finite witness spaces a computable brute-force extractor
does the same, so ``computable'' is no remedy. For R1CS this is a weakening, not a vacuity:
unsatisfiable instances exist, so RBR soundness remains a non-trivial property of the verifier, and
it is the property deployed STARKs are analysed for. Extractor \emph{efficiency} becomes
load-bearing in two places. One is the compiled knowledge claim, because the BCS extractor runs the
IOP extractor. The other is the recursion heuristic, because the CGSY24 bound contains
$N\cdot t_{\mathrm{ARG}}$. A contract that wants computational content must therefore demand a
named, executable extractor (the \lean{\ldots With} forms) and, at a higher level, a cost bound for
it.
